\section{Evaluasi}
\label{sec:evaluasi-bab1}

Evaluasi pada bab ini berbentuk \textbf{kuis diagnostik}. Kuis ini
tidak dinilai, tetapi wajib dikerjakan karena hasilnya menjadi
titik tolak bagi dosen dan mahasiswa dalam merencanakan pembelajaran
sepanjang semester.

\subsection{Soal Pilihan Ganda}

\begin{enumerate}
    \item Pada perancangan OBE, urutan yang benar adalah \dots
          \begin{enumerate}
              \item materi $\rightarrow$ asesmen $\rightarrow$ CPMK
                    $\rightarrow$ CPL
              \item CPL $\rightarrow$ CPMK $\rightarrow$ Sub-CPMK
                    $\rightarrow$ asesmen $\rightarrow$ materi
              \item CPMK $\rightarrow$ materi $\rightarrow$ CPL
                    $\rightarrow$ asesmen
              \item asesmen $\rightarrow$ materi $\rightarrow$ CPL
                    $\rightarrow$ CPMK
          \end{enumerate}

    \item Indikator yang lebih terukur dan dapat diuji, yang
          diturunkan dari CPMK, disebut \dots
          \begin{enumerate}
              \item CPL \item CPMK
              \item Sub-CPMK \item rubrik
          \end{enumerate}

    \item Bab yang membahas segmentasi memori dan alamat fisik adalah
          \dots
          \begin{enumerate}
              \item Bab~V \item Bab~VI
              \item Bab~VII \item Bab~VIII
          \end{enumerate}

    \item CPMK yang didukung oleh Bab~X sampai Bab~XIII adalah \dots
          \begin{enumerate}
              \item \CPMK{01} \item \CPMK{02}
              \item \CPMK{03} \item \CPMK{04}
          \end{enumerate}

    \item Komposisi penilaian mata kuliah ini adalah \dots
          \begin{enumerate}
              \item \dec{50}\% teori dan \dec{50}\% praktik
              \item \dec{30}\% teori dan \dec{70}\% praktik
              \item \dec{70}\% teori dan \dec{30}\% praktik
              \item \dec{100}\% ujian tertulis
          \end{enumerate}
\end{enumerate}

\subsection{Soal Uraian}

\begin{enumerate}
    \item Jelaskan mengapa pada OBE materi berada pada urutan terakhir
          dalam perancangan. Sertakan satu contoh materi yang mungkin
          \emph{tersisih} karena tidak mendukung indikator apa pun.

    \item Anda adalah mahasiswa yang baru masuk pada minggu ke-9.
          Bab mana yang harus Anda baca, dan Sub-CPMK apa yang sedang
          dibangun? Jelaskan alasan jawaban Anda.

    \item Bandingkan cara Anda mempersiapkan diri untuk komponen
          UTS (berbobot \dec{20}\%) dan komponen praktikum (berbobot
          \dec{20}\%). Apa perbedaan strategi belajarnya?

    \item Tuliskan rencana belajar Anda untuk empat minggu pertama
          dalam bentuk tabel: minggu, bab yang dibaca, Sub-CPMK yang
          ditargetkan, dan bukti belajar yang akan dikumpulkan.
\end{enumerate}

\subsection{Kisi-kisi Kuis Diagnostik}

\begin{table}[H]
    \centering
    \caption{Kisi-kisi kuis diagnostik Bab~\ref{chap:pendahuluan}.}
    \label{tab:kisi-bab1}
    \small
    \begin{tabular}{@{}p{5.6cm}p{5.0cm}c@{}}
        \toprule
        \textbf{Aspek yang diperiksa} & \textbf{Bentuk soal} &
        \textbf{Bobot} \\
        \midrule
        Pemahaman kerangka OBE & Pilihan ganda 1, 2; Uraian 1 &
        30\% \\
        \addlinespace
        Kemampuan menavigasi buku dan RPS & Pilihan ganda 3, 4;
        Uraian 2 & 30\% \\
        \addlinespace
        Pemahaman sistem penilaian & Pilihan ganda 5; Uraian 3 &
        20\% \\
        \addlinespace
        Kemampuan merencanakan belajar & Uraian 4 & 20\% \\
        \bottomrule
    \end{tabular}
\end{table}

\begin{checklistbox}
\textbf{Checklist Kesiapan Mengikuti Mata Kuliah}

Seorang mahasiswa dinyatakan siap mengikuti mata kuliah ini apabila:
\begin{itemize}
    \item[\checked] dapat menjelaskan tujuan buku dan keempat
          kemampuan akhirnya;
    \item[\checked] dapat menjelaskan jalur capaian pembelajaran dari
          CPL sampai asesmen;
    \item[\checked] dapat menemukan bab yang bersesuaian dengan
          minggu perkuliahan;
    \item[\checked] memiliki direktori kerja, jurnal praktikum, dan
          target belajar pribadi;
    \item[\checked] mengetahui komposisi penilaian dan letak rubrik
          penilaiannya.
\end{itemize}

\vspace{0.4em}
Bila ada butir yang belum bertanda, ulangi Bagian~\ref{sec:cara-menggunakan}
sebelum melanjutkan.
\end{checklistbox}
